\section{Introduction}
\label{sec:introduction}

Large Language Models (LLMs) have demonstrated strong performance across a wide
range of question-answering tasks. In the legal domain, however, a reliable
answer requires more than fluent generation: the system must identify relevant
provisions and reconstruct the reasoning chain connecting conditions,
intermediate events, and legal consequences
\cite{chalkidis2020legalbert,guha2023legalbench}.

RAG augments an LLM with evidence retrieved before generation
\cite{lewis2020retrieval,gao2024ragsurvey}. Conventional RAG is nevertheless
primarily driven by semantic similarity and often treats passages as
independent evidence. This is insufficient for multi-hop legal questions whose
conclusion depends on a sequence of connected condition--consequence rules.

We propose \textbf{LegalSCM-RAG}, which combines directed legal-graph retrieval
with question-aware structural verification. The verifier routes direct claims
to grounded edge checking, ordinary causal questions to three-valued factual
rule inference, and explicit mediator-removal questions to hard intervention.
A conservative commit gate uses a verifier decision only when grounding, path
coverage, confidence, and state-consistency requirements are satisfied;
otherwise it preserves the Causal Path RAG context and LLM decision.

The main contributions are:
\begin{itemize}
    \item a three-valued fixed-point formulation with an explicit sparse
    body/head incidence encoding, in which conjunction occurs within a rule and
    alternative disjunctive support across rules;
    \item a formal and implementation-level separation between graph
    node-deletion reachability and mechanism-replacement $do(X_M=0)$; and
    \item a controlled four-system evaluation with paired significance tests,
    abstention analysis, and route/commit auditing.
\end{itemize}

On the current benchmark, selective verification significantly improves
structured decision accuracy over Causal Path RAG, but not free-form answer or
citation quality. Importantly, the evaluated batch contains factual-validation
and direct-edge routes but no executed mediator intervention. We therefore
report the measured gain as evidence for the selective verifier as a whole,
not as an empirical validation of structural counterfactual intervention.
